\section{Synoptic Reference Tables}

The three tables gathered below place the core findings of Scripture, doctrine, and geography side by side for quick consultation during study or teaching.

\subsection{The Four Marian Dogmas and the Mystery of the Living Ark}

\begin{fourstable}{Dogma}{Defined Object and Authority}{Biblical Center and Reception}{Christological Coherence}
Divine Motherhood (\textit{Theotokos}) &
Mary truly conceived and bore the eternal Son of God incarnate in the unity of his divine Person. Confessed by Ephesus (431), Chalcedon (451), and Constantinople II (553, Canon VI). &
Luke 1:31--35, 43 (``the mother of my Lord''); Gal 4:4. Witnessed early by Athanasius, Cyril of Jerusalem, Gregory of Nazianzus, Cyril of Alexandria, John of Damascus, and Aquinas (\textit{ST} III, q. 35, a. 4). &
As the old Ark was holy solely through the Lord enthroned above it, Mary is Mother of God solely because of the divine Son who took flesh from her. \\
\addlinespace[0.35em]
Perpetual Virginity (\textit{Aeiparthenos}) &
Mary remained a virgin before, during, and after the birth of Christ. Named \textit{ever-virgin} by Constantinople II (553, Canons VI, XIV); articulated in threefold form at Lateran (649, Canon 3) and CCC 499. &
Isa 7:14; Matt 1:18--25; Luke 1:27, 34--35. Defended by Jerome (\textit{Contra Helvidium}), Ambrose (\textit{In Lucam} II), Augustine (\textit{Sermon} 51.18), Bernard, and Aquinas (\textit{ST} III, q. 28). &
Her virginal body and soul are set apart without division as the shrine of the Holy Spirit (\textit{Sacrarium Spiritus Sancti}) and icon of the spotless Virgin-Mother Church. \\
\addlinespace[0.35em]
Immaculate Conception &
Preserved from all stain of original sin from the first instant of her conception, by singular grace in view of Christ's foreseen merits. Defined \textit{ex cathedra} by Pius IX, \textit{Ineffabilis Deus} (8 Dec 1854). &
Gen 3:15; Luke 1:28 (\textit{kecharitomene}); Eph 1:3--6. Augustine (\textit{De natura et gratia} 36); Aquinas's hesitation (\textit{ST} III, q. 27, a. 2) resolved via Duns Scotus's preservative redemption; Trent (Sess. V, \S 6). &
Prefigured by pure gold, incorruptible acacia wood, and the Tabernacle crafted by the new Bezalel (\textit{Ineffabilis Deus}); Christ's foreseen Passion prepared his own spotless dwelling. \\
\addlinespace[0.35em]
Bodily Assumption &
Having completed the course of her earthly life, Mary was assumed body and soul into heavenly glory. Defined \textit{ex cathedra} by Pius XII, \textit{Munificentissimus Deus} (1 Nov 1950, \S 44). Passing through death remains open. &
Gen 3:15; Ps 131(132):8; 1 Cor 15:20--27; Apoc 11:19--12:1. Organic growth from Augustine's ecclesial reading of Ps 131:8 through John Damascene's Dormition homilies, Aquinas (\textit{ST} III, q. 27, a. 1), and Assumption liturgy. &
The living Ark that carried the Author of life is not abandoned to corruption; she enters ahead of the pilgrim Church into the resting place opened by her Son's Resurrection. \\
\end{fourstable}

\subsection{Exodus 40, 2 Samuel 6, and Luke 1: Graded Comparison}

\begin{fourstable}{Motif}{Old Testament Passage}{Luke 1 Passage}{Level of Correspondence}
Divine overshadowing &
Brenton LXX Exod 40:35 (``because the cloud overshadowed it''); Num 9:18--22; 10:36; Swete LXX Ps 90:4 (\textit{episkiasei soi}; Vulgate \textit{obumbrabit tibi}). &
Luke 1:35: ``the power of the Most High shall overshadow thee'' (\textit{episkiasei soi}; Vulgate \textit{obumbrabit tibi}). One of 5 NT tokens of \textit{episkiazo}. &
Direct Septuagint--New Testament lexical link (exact inflected match at Ps 90:4; Vulgate Exodus 40:33 uses \textit{operio}, not \textit{obumbro}). \\
\addlinespace[0.35em]
Arising and going into Judah &
2 Sam 6:2: David arose and went into Judah (MT \textit{Baale of Judah} = Kiriath-jearim; Vulgate/DR \textit{de viris Juda}). &
Luke 1:39: Mary ``rising up ... went into the hill country with haste into a city of Juda'' (town unnamed; trad. Ein Karem). &
Close narrative and geographic correspondence; identical Hebrew-Greek idiom of arising to go, without claiming an identical road. \\
\addlinespace[0.35em]
Liturgical outcry &
1 Chr 15:16, 28; 16:4--5, 42; 2 Chr 5:13 (Brenton LXX Levitical shouting and lifting up the voice before the Ark); 2 Sam 6:15. &
Luke 1:42: Elizabeth cried out (\textit{anephonesen}, exact NT hapax legomenon; SBLGNT \textit{krauge megale}, RP2018 \textit{phone megale}). &
Unique New Testament verb aligned with the Chronicler's Levitical Ark-proclamation idiom. \\
\addlinespace[0.35em]
Reverent question at arrival &
2 Sam 6:9: ``How shall the ark of the Lord come to me?'' (\textit{Quomodo ingredietur ad me arca Domini?}); 1 Chr 13:12. &
Luke 1:43: ``And whence is this to me that the mother of my Lord should come to me?'' (\textit{Et unde hoc mihi, ut veniat mater Domini mei ad me?}). &
Tight structural and syntactic parallel expressing awe at the coming of the Lord's bearer, with distinct interrogatives (\textit{pos} vs. \textit{pothen}). \\
\addlinespace[0.35em]
Bodily joy before the Lord &
2 Sam 6:14, 16: David danced / leaping and dancing (Vulgate \textit{saltabat}, \textit{subsilientem}; Brenton LXX ``sounded with well-tuned instruments'' / ``dancing and playing''). &
Luke 1:41, 44: the infant leaped (\textit{eskirtesen}, Vulgate \textit{exsultavit}). \textit{Skirtao} occurs 3 times in NT, all in Luke; echoes Swete Ps 113:4, 6 and Gen 25:22. &
Thematic and narrative parallel expressed through different Greek and Latin verbs; never a shared verb between 2 Sam 6 and Luke 1. \\
\addlinespace[0.35em]
Three months of household blessing &
2 Sam 6:11; 1 Chr 13:14: the Ark remained three months in Obed-edom's house (\textit{Mansit ... tribus mensibus}) and the Lord blessed all his household. &
Luke 1:56: ``And Mary abode with her about three months'' (\textit{Mansit autem Maria cum illa quasi mensibus tribus}) in the priestly house of Zachary and Elizabeth. &
Direct narrative-temporal parallel with identical Vulgate cadence (\textit{mansit ... tribus mensibus}); note patristic commentaries (Ambrose, Lapide) did not yet draw the Obed-edom comparison at 1:56. \\
\end{fourstable}

\subsection{Chronological Arc of the Ark's Canonical Stations}

\begin{threestable}{Stage}{Scripture}{Controlling Theological Significance}
Sinai and the Wilderness &
Exod 25, 37, 40; Num 4, 7, 10; Deut 10, 31 &
Commanded from the innermost heart outward; overshadowed by the glory-cloud; borne on Kohathite shoulders as a pilgrim throne going ahead of Israel. \\
\addlinespace[0.3em]
Jordan, Jericho, Ebal, and Shiloh &
Jos 3--4, 6, 8:30--35, 18:1; Judg 20:18, 26--28; 1 Sam 1--3 &
Parts the floodwaters, casts down walls by liturgical obedience, witnesses covenant renewal, and settles at Shiloh (noting the Bethel reading in MT/LXX Judges 20). \\
\addlinespace[0.3em]
Philistine Captivity and Hiddenness &
1 Sam 4--7 &
Refuses to act as a battlefield talisman; shatters Dagon in Ashdod; moves through Gath and Ekron (Brenton LXX Ashkelon at 5:10); rests twenty hidden years at Kiriath-jearim. \\
\addlinespace[0.3em]
Perez-Uzzah, Obed-edom, and Zion &
2 Sam 6; 1 Chr 13, 15--16; Ps 131(132):8 &
Judges disordered handling, pours blessing on Obed-edom's reverent household for three months, and ascends amid Levitical praise and David's dance to Mount Zion. \\
\addlinespace[0.3em]
Solomon's Temple, Silence, and Prophecy &
1 Kgs 6--8; 2 Chr 5, 35:3; Jer 3:16; 2 Macc 2:1--8 &
Enters the cedar \textit{debir} with only the two tablets inside and poles still visible; last commanded in Josiah's reign; hidden in exile while prophecy points beyond a wooden chest. \\
\addlinespace[0.3em]
Incarnation, Visitation, and Heavenly Temple &
Luke 1--2; Matt 1--2; Heb 9:1--12; Apoc 11:19--12:6 &
Fulfilled in the Word made flesh within Mary's overshadowed womb, borne in haste into Judah, guarded by Joseph, and unveiled in the opened temple of heaven beside the Woman clothed with the sun. \\
\end{threestable}
